% =====================================================================
%  សន្លឹកកិច្ចការ — រូបវិទ្យា ថ្នាក់ទី០៩ · មេរៀនទី០១ ៖ ការបំបែកកម្លាំង
%  រៀបចំដោយ៖ ពៅ គឹមយ៉ាត — Science Master Academic
%
%  Compile with XeLaTeX.
%    Student copy : xelatex physics-grade9-lesson01-worksheet.tex
%    Teacher copy : xelatex -jobname=physics-grade9-lesson01-worksheet-answers \
%                     "\def\showanswers{}% =====================================================================
%  សន្លឹកកិច្ចការ — រូបវិទ្យា ថ្នាក់ទី០៩ · មេរៀនទី០១ ៖ ការបំបែកកម្លាំង
%  រៀបចំដោយ៖ ពៅ គឹមយ៉ាត — Science Master Academic
%
%  Compile with XeLaTeX.
%    Student copy : xelatex physics-grade9-lesson01-worksheet.tex
%    Teacher copy : xelatex -jobname=physics-grade9-lesson01-worksheet-answers \
%                     "\def\showanswers{}% =====================================================================
%  សន្លឹកកិច្ចការ — រូបវិទ្យា ថ្នាក់ទី០៩ · មេរៀនទី០១ ៖ ការបំបែកកម្លាំង
%  រៀបចំដោយ៖ ពៅ គឹមយ៉ាត — Science Master Academic
%
%  Compile with XeLaTeX.
%    Student copy : xelatex physics-grade9-lesson01-worksheet.tex
%    Teacher copy : xelatex -jobname=physics-grade9-lesson01-worksheet-answers \
%                     "\def\showanswers{}% =====================================================================
%  សន្លឹកកិច្ចការ — រូបវិទ្យា ថ្នាក់ទី០៩ · មេរៀនទី០១ ៖ ការបំបែកកម្លាំង
%  រៀបចំដោយ៖ ពៅ គឹមយ៉ាត — Science Master Academic
%
%  Compile with XeLaTeX.
%    Student copy : xelatex physics-grade9-lesson01-worksheet.tex
%    Teacher copy : xelatex -jobname=physics-grade9-lesson01-worksheet-answers \
%                     "\def\showanswers{}\input{physics-grade9-lesson01-worksheet}"
%  Requires the Khmer font "Noto Sans Khmer".
% =====================================================================
\documentclass[11pt,a4paper]{article}

\usepackage{amsmath,amssymb}
\usepackage{fontspec}
\usepackage[a4paper,top=1.4cm,bottom=1.9cm,left=1.8cm,right=1.8cm,
  footskip=0.9cm]{geometry}
\usepackage{xcolor}
\usepackage{tikz}
\usetikzlibrary{calc}
\usepackage[most]{tcolorbox}
\usepackage{fancyhdr}

% ---- Fonts: Latin = Latin Modern Sans, Khmer = Noto Sans Khmer ----
\setmainfont{Latin Modern Sans}
\newfontfamily\khmerfont[Script=Khmer]{Noto Sans Khmer}
\usepackage[Latin,Khmer]{ucharclasses}
\setTransitionsFor{Khmer}{\khmerfont}{\rmfamily}
\linespread{1.25}
\setlength{\parindent}{0pt}

% ---- Colours ----
\definecolor{brand}{HTML}{1F6B5E}
\definecolor{mint}{HTML}{EEF6F3}
\definecolor{ink}{HTML}{1E2A2A}
\definecolor{inkgray}{HTML}{5D6B6B}
\definecolor{linegray}{HTML}{C9D6D3}
\definecolor{rightfill}{HTML}{DDF1E6}
\definecolor{crimson}{HTML}{B3263A}

% ---- Student / teacher mode ----
\newif\ifanswers
\ifdefined\showanswers \answerstrue \fi

\def\W{17.4} % text width in cm

% ---- Khmer numerals ----
\newcommand{\khnum}[1]{\ifcase#1\relax ០\or ១\or ២\or ៣\or ៤\or ៥\or ៦\or ៧\or
  ៨\or ៩\or ១០\or ១១\or ១២\or ១៣\or ១៤\or ១៥\or ១៦\or ១៧\or ១៨\or ១៩\or ២០\fi}
\newcommand{\khletter}[1]{\ifcase#1\relax\or ក\or ខ\or គ\or ឃ\fi}

% ---- Footer ----
\pagestyle{fancy}
\fancyhf{}
\renewcommand{\headrulewidth}{0pt}
\renewcommand{\footrulewidth}{0.5pt}
\renewcommand{\footrule}{{\color{linegray}\hrule height\footrulewidth width\headwidth}\vspace{3pt}}
\fancyfoot[L]{\scriptsize\color{inkgray}Science Master Academic · រូបវិទ្យា ថ្នាក់ទី០៩ ·
  ទំព័រ \khnum{\value{page}} / ២០}
\fancyfoot[R]{\scriptsize\color{inkgray}ទំនាក់ទំនង 088 909 3291 · Telegram @POVKIMYAT}

% ---- Page header ----
\newcommand{\pageheader}{%
\begin{tikzpicture}
  \useasboundingbox (0,0.1) rectangle (\W,-2.45);
  % logo
  \begin{scope}[shift={(0.9,-1.05)}]
    \fill[mint] (0,0) circle (0.88);
    \draw[brand,line width=1.1pt] (0,0) circle (0.88);
    \foreach \x/\y in {-0.19/0.36,0.19/0.36,-0.19/-0.02,0.19/-0.02}
      \draw[brand,line width=1.1pt,rounded corners=1.5pt]
        (\x-0.15,\y-0.15) rectangle (\x+0.15,\y+0.15);
    \node[font=\fontsize{5}{5}\selectfont\bfseries,text=brand] at (0,-0.5) {Science Master};
  \end{scope}
  % titles
  \node[anchor=west,font=\small\bfseries,text=brand] at (2.05,-0.3) {SCIENCE MASTER ACADEMIC};
  \node[anchor=west,font=\LARGE\bfseries,text=ink] at (2.05,-1.0)
    {សន្លឹកកិច្ចការ\ifanswers{\large\color{crimson}\ (ចម្លើយសម្រាប់គ្រូ)}\fi};
  \node[anchor=west,font=\small\bfseries,text=ink] at (2.05,-1.75)
    {រូបវិទ្យា ថ្នាក់ទី០៩ · មេរៀនទី០១ ៖ ការបំបែកកម្លាំង};
  % score box
  \draw[linegray,rounded corners=6pt,line width=0.9pt] (\W-2.6,-0.15) rectangle (\W,-1.75);
  \node[font=\footnotesize,text=inkgray] at (\W-1.3,-0.45) {ពិន្ទុ};
  \node[font=\large,text=ink] at (\W-1.3,-1.2) {\rule{1cm}{0.5pt}\, / 10};
  % rule
  \draw[brand,line width=2pt] (0,-2.35) -- (\W,-2.35);
\end{tikzpicture}\par\vspace{8pt}
{\small ឈ្មោះសិស្ស ៖ \makebox[6.2cm]{\color{inkgray}\dotfill}\hfill
 កាលបរិច្ឆេទ ៖ \makebox[4.2cm]{\color{inkgray}\dotfill}}\par\vspace{8pt}}

% ---- Difficulty badge ----
\newcommand{\badge}[1]{\tcbox[on line,colback=white,colframe=brand,boxrule=0.7pt,
  arc=3pt,left=4pt,right=4pt,top=0pt,bottom=0pt,boxsep=1pt]{\scriptsize\bfseries\color{brand}#1}}

% ---- Multiple-choice options ----
\tcbset{optbase/.style={enhanced,arc=6pt,boxrule=0.7pt,colframe=linegray,
  left=6pt,right=6pt,top=3pt,bottom=3pt,valign=center}}
\newcommand{\optlabel}[1]{\tikz[baseline=(c.base)]\node[draw=brand,circle,line width=0.8pt,
  inner sep=1.2pt,minimum size=1.35em,font=\footnotesize\bfseries,text=ink](c){\khletter{#1}};}
\newcommand{\opt}[2]{%
  \ifanswers\ifnum#1=\correct\relax
    \begin{tcolorbox}[optbase,colback=rightfill,colframe=brand,boxrule=1.2pt]
      \optlabel{#1}\quad #2\hfill{\color{brand}$\checkmark$}\end{tcolorbox}
  \else
    \begin{tcolorbox}[optbase,colback=white]\optlabel{#1}\quad #2\end{tcolorbox}
  \fi\else
    \begin{tcolorbox}[optbase,colback=white]\optlabel{#1}\quad #2\end{tcolorbox}
  \fi}

% ---- Ruled writing lines (student copy only) ----
\ifanswers
  \tcbset{writelines/.style={}}
\else
  \tcbset{writelines/.style={underlay={\begin{tcbclipinterior}
      \foreach \i in {1,...,30}
        \draw[linegray,line width=0.4pt]
          ($(interior.north west)+(0.5,-\i*0.72+0.05)$) --
          ($(interior.north east)+(-0.5,-\i*0.72+0.05)$);
      \end{tcbclipinterior}}}}
\fi

% ---- Solution text (shown only in teacher copy) ----
\newcommand{\sol}[1]{\gdef\thesol{#1}}

% ---- One exercise = one page ----
\newcounter{prob}
% #1 level  #2 intro  #3 display  #4 question  #5-#8 options  #9 correct (1-4)
\newcommand{\problem}[9]{%
  \stepcounter{prob}\def\correct{#9}%
  \pageheader
  \begin{tcolorbox}[enhanced,colback=mint,colframe=brand,boxrule=0.9pt,arc=8pt,
    left=10pt,right=10pt,top=8pt,bottom=8pt]
    {\bfseries\color{ink}លំហាត់ទី \khnum{\value{prob}}}\enspace\badge{#1}\par\vspace{2pt}
    {\small #2}\par\vspace{6pt}
    \begin{tcolorbox}[enhanced,colback=white,colframe=linegray,boxrule=0.6pt,arc=6pt,
      halign=center,top=6pt,bottom=6pt]
      {\Large\mbox{$\displaystyle #3$}}\par\vspace{4pt}
      {\small\bfseries #4}
    \end{tcolorbox}
  \end{tcolorbox}
  \vspace{6pt}
  {\bfseries ក. គូសរង្វង់ចម្លើយដែលត្រឹមត្រូវ}\par\vspace{4pt}
  \begin{tcbraster}[raster columns=2,raster equal height,raster column skip=12pt,
    raster row skip=7pt]
    \opt{1}{#5}\opt{2}{#6}\opt{3}{#7}\opt{4}{#8}
  \end{tcbraster}
  \vspace{8pt}
  {\bfseries ខ. សរសេរដំណោះស្រាយលម្អិត}\par\vspace{4pt}
  \begin{tcolorbox}[enhanced,height fill,colback=white,colframe=linegray,boxrule=0.6pt,
    arc=8pt,left=14pt,right=14pt,top=10pt,writelines]
    \ifanswers
      {\color{crimson}\bfseries ចម្លើយ ៖ (\khletter{#9})}\par\vspace{2pt}
      {\color{crimson}\thesol}
    \else\mbox{}\fi
  \end{tcolorbox}
  \newpage}

\newcommand{\N}{\,\mathrm{N}}
\newcommand{\kg}{\,\mathrm{kg}}

\begin{document}

% ================= 1 =================
\sol{$F_x=F\cos\alpha=10\times\cos60^\circ=10\times0.5=5\N$}
\problem{ងាយ}
  {កម្លាំង $\vec F$ ធ្វើមុំ $\alpha$ ជាមួយអ័ក្សផ្ដេក $Ox$ ។ គេឲ្យ}
  {F=10\N\,,\qquad \alpha=60^\circ}
  {គណនាកម្លាំងបំបែកផ្ដេក $F_x$}
  {$5\N$}{$8.66\N$}{$10\N$}{$0\N$}{1}

% ================= 2 =================
\sol{$F_y=F\sin\alpha=10\times\sin60^\circ=10\times0.866\approx8.66\N$}
\problem{ងាយ}
  {កម្លាំង $\vec F$ ធ្វើមុំ $\alpha$ ជាមួយអ័ក្សផ្ដេក $Ox$ ។ គេឲ្យ}
  {F=10\N\,,\qquad \alpha=60^\circ}
  {គណនាកម្លាំងបំបែកបញ្ឈរ $F_y$}
  {$5\N$}{$10\N$}{$8.66\N$}{$17.3\N$}{3}

% ================= 3 =================
\sol{$F_y=F\sin\alpha=20\times\sin30^\circ=20\times0.5=10\N$}
\problem{ងាយ}
  {កម្លាំង $\vec F$ ធ្វើមុំ $\alpha$ ជាមួយអ័ក្សផ្ដេក $Ox$ ។ គេឲ្យ}
  {F=20\N\,,\qquad \alpha=30^\circ}
  {គណនាកម្លាំងបំបែកបញ្ឈរ $F_y$}
  {$17.3\N$}{$10\N$}{$20\N$}{$40\N$}{2}

% ================= 4 =================
\sol{$F_x=F\cos45^\circ=20\times0.707\approx14.1\N$}
\problem{ងាយ}
  {កម្លាំង $\vec F$ ធ្វើមុំ $\alpha$ ជាមួយអ័ក្សផ្ដេក $Ox$ ។ គេឲ្យ}
  {F=20\N\,,\qquad \alpha=45^\circ}
  {គណនាកម្លាំងបំបែកផ្ដេក $F_x$}
  {$10\N$}{$20\N$}{$28.3\N$}{$14.1\N$}{4}

% ================= 5 =================
\sol{$F_x=F\cos37^\circ=50\times0.8=40\N$}
\problem{ងាយ}
  {កម្លាំង $\vec F$ ធ្វើមុំ $\alpha$ ជាមួយអ័ក្សផ្ដេក ។ (យក $\sin37^\circ=0.6$ ,
   $\cos37^\circ=0.8$) ។ គេឲ្យ}
  {F=50\N\,,\qquad \alpha=37^\circ}
  {គណនាកម្លាំងបំបែកផ្ដេក $F_x$}
  {$40\N$}{$30\N$}{$50\N$}{$25\N$}{1}

% ================= 6 =================
\sol{$F_y=F\sin37^\circ=50\times0.6=30\N$}
\problem{ងាយ}
  {កម្លាំង $\vec F$ ធ្វើមុំ $\alpha$ ជាមួយអ័ក្សផ្ដេក ។ (យក $\sin37^\circ=0.6$ ,
   $\cos37^\circ=0.8$) ។ គេឲ្យ}
  {F=50\N\,,\qquad \alpha=37^\circ}
  {គណនាកម្លាំងបំបែកបញ្ឈរ $F_y$}
  {$40\N$}{$20\N$}{$30\N$}{$50\N$}{3}

% ================= 7 =================
\sol{$F=\sqrt{F_x^2+F_y^2}=\sqrt{3^2+4^2}=\sqrt{25}=5\N$}
\problem{ងាយ}
  {កម្លាំង $\vec F$ ត្រូវបានបំបែកជាកម្លាំងពីរកែងគ្នា ។ គេឲ្យ}
  {F_x=3\N\,,\qquad F_y=4\N}
  {គណនាអាំងតង់ស៊ីតេនៃកម្លាំងផ្គុំ $F$}
  {$7\N$}{$5\N$}{$1\N$}{$12\N$}{2}

% ================= 8 =================
\sol{$F=\sqrt{6^2+8^2}=\sqrt{100}=10\N$}
\problem{ងាយ}
  {កម្លាំង $\vec F$ ត្រូវបានបំបែកជាកម្លាំងពីរកែងគ្នា ។ គេឲ្យ}
  {F_x=6\N\,,\qquad F_y=8\N}
  {គណនាអាំងតង់ស៊ីតេនៃកម្លាំងផ្គុំ $F$}
  {$14\N$}{$2\N$}{$48\N$}{$10\N$}{4}

% ================= 9 =================
\sol{$F=\sqrt{12^2+5^2}=\sqrt{144+25}=\sqrt{169}=13\N$}
\problem{មធ្យម}
  {កម្លាំង $\vec F$ ត្រូវបានបំបែកជាកម្លាំងពីរកែងគ្នា ។ គេឲ្យ}
  {F_x=12\N\,,\qquad F_y=5\N}
  {គណនាអាំងតង់ស៊ីតេនៃកម្លាំងផ្គុំ $F$}
  {$17\N$}{$7\N$}{$13\N$}{$60\N$}{3}

% ================= 10 =================
\sol{$\tan\alpha=\dfrac{F_y}{F_x}=\dfrac{10}{10}=1\;\Rightarrow\;\alpha=45^\circ$}
\problem{មធ្យម}
  {កម្លាំង $\vec F$ មានកម្លាំងបំបែកកែងគ្នាពីរ ។ គេឲ្យ}
  {F_x=10\N\,,\qquad F_y=10\N}
  {រកមុំ $\alpha$ ដែល $\vec F$ ធ្វើជាមួយអ័ក្សផ្ដេក $Ox$}
  {$45^\circ$}{$30^\circ$}{$60^\circ$}{$90^\circ$}{1}

% ================= 11 =================
\sol{$\tan\alpha=\dfrac{F_y}{F_x}=\dfrac{10}{10\sqrt3}=\dfrac{1}{\sqrt3}\;\Rightarrow\;\alpha=30^\circ$}
\problem{មធ្យម}
  {កម្លាំង $\vec F$ មានកម្លាំងបំបែកកែងគ្នាពីរ ។ គេឲ្យ}
  {F_x=10\sqrt3\N\,,\qquad F_y=10\N}
  {រកមុំ $\alpha$ ដែល $\vec F$ ធ្វើជាមួយអ័ក្សផ្ដេក $Ox$}
  {$60^\circ$}{$30^\circ$}{$45^\circ$}{$90^\circ$}{2}

% ================= 12 =================
\sol{$\cos\alpha=\dfrac{F_x}{F}=\dfrac{20}{40}=0.5\;\Rightarrow\;\alpha=60^\circ$}
\problem{មធ្យម}
  {កម្លាំង $\vec F$ ធ្វើមុំ $\alpha$ ជាមួយអ័ក្សផ្ដេក $Ox$ ។ គេឲ្យ}
  {F=40\N\,,\qquad F_x=20\N}
  {រកមុំ $\alpha$}
  {$30^\circ$}{$45^\circ$}{$90^\circ$}{$60^\circ$}{4}

% ================= 13 =================
\sol{$F_x=F\cos\alpha\;\Rightarrow\;F=\dfrac{F_x}{\cos60^\circ}=\dfrac{15}{0.5}=30\N$}
\problem{មធ្យម}
  {កម្លាំង $\vec F$ ធ្វើមុំ $\alpha$ ជាមួយអ័ក្សផ្ដេក $Ox$ ។ គេឲ្យ}
  {F_x=15\N\,,\qquad \alpha=60^\circ}
  {រកអាំងតង់ស៊ីតេនៃកម្លាំង $F$}
  {$7.5\N$}{$15\N$}{$30\N$}{$26\N$}{3}

% ================= 14 =================
\sol{មុំជាមួយអ័ក្សផ្ដេក $\alpha=90^\circ-60^\circ=30^\circ$ ។\par
     $F_x=F\cos30^\circ=30\times0.866\approx26\N$}
\problem{មធ្យម}
  {កម្លាំង $\vec F$ ធ្វើមុំ $\beta$ ជាមួយអ័ក្ស\textbf{បញ្ឈរ} $Oy$ ។ គេឲ្យ}
  {F=30\N\,,\qquad \beta=60^\circ}
  {គណនាកម្លាំងបំបែកផ្ដេក $F_x$}
  {$26\N$}{$15\N$}{$30\N$}{$52\N$}{1}

% ================= 15 =================
\sol{កម្លាំងទាញទៅមុខ គឺកម្លាំងបំបែកផ្ដេក ៖\par
     $F_x=F\cos60^\circ=40\times0.5=20\N$}
\problem{មធ្យម}
  {សិស្សម្នាក់ទាញរទេះដោយខ្សែ ដែលធ្វើមុំ $\alpha$ ជាមួយផ្ទៃដីផ្ដេក ។ គេឲ្យ}
  {F=40\N\,,\qquad \alpha=60^\circ}
  {រកកម្លាំងដែលទាញរទេះទៅមុខ (កម្លាំងបំបែកផ្ដេក)}
  {$34.6\N$}{$40\N$}{$80\N$}{$20\N$}{4}

% ================= 16 =================
\sol{$P=mg=5\times10=50\N$\par
     $P_{\parallel}=P\sin30^\circ=50\times0.5=25\N$}
\problem{មធ្យម}
  {ប្រអប់ម៉ាស $m$ ស្ថិតលើប្លង់ទេរ ដែលមានមុំទេរ $\alpha$ ។ ($g=10\,\mathrm{N/kg}$) ។ គេឲ្យ}
  {m=5\kg\,,\qquad \alpha=30^\circ}
  {គណនាកម្លាំងបំបែកនៃទម្ងន់ ស្របនឹងប្លង់ទេរ $P_{\parallel}$}
  {$43.3\N$}{$25\N$}{$50\N$}{$5\N$}{2}

% ================= 17 =================
\sol{$P=mg=5\times10=50\N$\par
     $P_{\perp}=P\cos30^\circ=50\times0.866\approx43.3\N$}
\problem{មធ្យម}
  {ប្រអប់ម៉ាស $m$ ស្ថិតលើប្លង់ទេរ ដែលមានមុំទេរ $\alpha$ ។ ($g=10\,\mathrm{N/kg}$) ។ គេឲ្យ}
  {m=5\kg\,,\qquad \alpha=30^\circ}
  {គណនាកម្លាំងបំបែកនៃទម្ងន់ កែងនឹងប្លង់ទេរ $P_{\perp}$}
  {$25\N$}{$50\N$}{$43.3\N$}{$86.6\N$}{3}

% ================= 18 =================
\sol{$P=mg=2\times10=20\N$\par
     $P_{\parallel}=P\sin37^\circ=20\times0.6=12\N$}
\problem{មធ្យម}
  {ប្រអប់ម៉ាស $m$ ស្ថិតលើប្លង់ទេរ ដែលមានមុំទេរ $\alpha$ ។ ($g=10\,\mathrm{N/kg}$ ,
   $\sin37^\circ=0.6$ , $\cos37^\circ=0.8$) ។ គេឲ្យ}
  {m=2\kg\,,\qquad \alpha=37^\circ}
  {គណនាកម្លាំងបំបែកនៃទម្ងន់ ស្របនឹងប្លង់ទេរ $P_{\parallel}$}
  {$12\N$}{$16\N$}{$20\N$}{$6\N$}{1}

% ================= 19 =================
\sol{$P_{\parallel}$ និង $P_{\perp}$ កែងគ្នា ដូច្នេះ\par
     $P=\sqrt{P_{\parallel}^2+P_{\perp}^2}=\sqrt{30^2+40^2}=\sqrt{2500}=50\N$}
\problem{ពិបាក}
  {ទម្ងន់ $\vec P$ របស់វត្ថុមួយលើប្លង់ទេរ ត្រូវបានបំបែកជាកម្លាំងពីរ ។ គេឲ្យ}
  {P_{\parallel}=30\N\,,\qquad P_{\perp}=40\N}
  {រកទម្ងន់ $P$ របស់វត្ថុ}
  {$70\N$}{$10\N$}{$35\N$}{$50\N$}{4}

% ================= 20 =================
\sol{ដើម្បីទប់ប្រអប់ $F$ ត្រូវស្មើនឹង $P_{\parallel}$ ៖\par
     $F=P_{\parallel}=mg\sin30^\circ=4\times10\times0.5=20\N$}
\problem{ពិបាក}
  {ប្រអប់ម៉ាស $m$ ស្ថិតលើប្លង់ទេររលោង (គ្មានកកិត) ដែលមានមុំទេរ $\alpha$ ។
   ដើម្បីទប់ប្រអប់កុំឲ្យរអិល គេអនុវត្តកម្លាំង $\vec F$ ស្របនឹងប្លង់ ។
   ($g=10\,\mathrm{N/kg}$) ។ គេឲ្យ}
  {m=4\kg\,,\qquad \alpha=30^\circ}
  {រកអាំងតង់ស៊ីតេនៃកម្លាំង $F$}
  {$40\N$}{$20\N$}{$34.6\N$}{$10\N$}{2}

\end{document}
"
%  Requires the Khmer font "Noto Sans Khmer".
% =====================================================================
\documentclass[11pt,a4paper]{article}

\usepackage{amsmath,amssymb}
\usepackage{fontspec}
\usepackage[a4paper,top=1.4cm,bottom=1.9cm,left=1.8cm,right=1.8cm,
  footskip=0.9cm]{geometry}
\usepackage{xcolor}
\usepackage{tikz}
\usetikzlibrary{calc}
\usepackage[most]{tcolorbox}
\usepackage{fancyhdr}

% ---- Fonts: Latin = Latin Modern Sans, Khmer = Noto Sans Khmer ----
\setmainfont{Latin Modern Sans}
\newfontfamily\khmerfont[Script=Khmer]{Noto Sans Khmer}
\usepackage[Latin,Khmer]{ucharclasses}
\setTransitionsFor{Khmer}{\khmerfont}{\rmfamily}
\linespread{1.25}
\setlength{\parindent}{0pt}

% ---- Colours ----
\definecolor{brand}{HTML}{1F6B5E}
\definecolor{mint}{HTML}{EEF6F3}
\definecolor{ink}{HTML}{1E2A2A}
\definecolor{inkgray}{HTML}{5D6B6B}
\definecolor{linegray}{HTML}{C9D6D3}
\definecolor{rightfill}{HTML}{DDF1E6}
\definecolor{crimson}{HTML}{B3263A}

% ---- Student / teacher mode ----
\newif\ifanswers
\ifdefined\showanswers \answerstrue \fi

\def\W{17.4} % text width in cm

% ---- Khmer numerals ----
\newcommand{\khnum}[1]{\ifcase#1\relax ០\or ១\or ២\or ៣\or ៤\or ៥\or ៦\or ៧\or
  ៨\or ៩\or ១០\or ១១\or ១២\or ១៣\or ១៤\or ១៥\or ១៦\or ១៧\or ១៨\or ១៩\or ២០\fi}
\newcommand{\khletter}[1]{\ifcase#1\relax\or ក\or ខ\or គ\or ឃ\fi}

% ---- Footer ----
\pagestyle{fancy}
\fancyhf{}
\renewcommand{\headrulewidth}{0pt}
\renewcommand{\footrulewidth}{0.5pt}
\renewcommand{\footrule}{{\color{linegray}\hrule height\footrulewidth width\headwidth}\vspace{3pt}}
\fancyfoot[L]{\scriptsize\color{inkgray}Science Master Academic · រូបវិទ្យា ថ្នាក់ទី០៩ ·
  ទំព័រ \khnum{\value{page}} / ២០}
\fancyfoot[R]{\scriptsize\color{inkgray}ទំនាក់ទំនង 088 909 3291 · Telegram @POVKIMYAT}

% ---- Page header ----
\newcommand{\pageheader}{%
\begin{tikzpicture}
  \useasboundingbox (0,0.1) rectangle (\W,-2.45);
  % logo
  \begin{scope}[shift={(0.9,-1.05)}]
    \fill[mint] (0,0) circle (0.88);
    \draw[brand,line width=1.1pt] (0,0) circle (0.88);
    \foreach \x/\y in {-0.19/0.36,0.19/0.36,-0.19/-0.02,0.19/-0.02}
      \draw[brand,line width=1.1pt,rounded corners=1.5pt]
        (\x-0.15,\y-0.15) rectangle (\x+0.15,\y+0.15);
    \node[font=\fontsize{5}{5}\selectfont\bfseries,text=brand] at (0,-0.5) {Science Master};
  \end{scope}
  % titles
  \node[anchor=west,font=\small\bfseries,text=brand] at (2.05,-0.3) {SCIENCE MASTER ACADEMIC};
  \node[anchor=west,font=\LARGE\bfseries,text=ink] at (2.05,-1.0)
    {សន្លឹកកិច្ចការ\ifanswers{\large\color{crimson}\ (ចម្លើយសម្រាប់គ្រូ)}\fi};
  \node[anchor=west,font=\small\bfseries,text=ink] at (2.05,-1.75)
    {រូបវិទ្យា ថ្នាក់ទី០៩ · មេរៀនទី០១ ៖ ការបំបែកកម្លាំង};
  % score box
  \draw[linegray,rounded corners=6pt,line width=0.9pt] (\W-2.6,-0.15) rectangle (\W,-1.75);
  \node[font=\footnotesize,text=inkgray] at (\W-1.3,-0.45) {ពិន្ទុ};
  \node[font=\large,text=ink] at (\W-1.3,-1.2) {\rule{1cm}{0.5pt}\, / 10};
  % rule
  \draw[brand,line width=2pt] (0,-2.35) -- (\W,-2.35);
\end{tikzpicture}\par\vspace{8pt}
{\small ឈ្មោះសិស្ស ៖ \makebox[6.2cm]{\color{inkgray}\dotfill}\hfill
 កាលបរិច្ឆេទ ៖ \makebox[4.2cm]{\color{inkgray}\dotfill}}\par\vspace{8pt}}

% ---- Difficulty badge ----
\newcommand{\badge}[1]{\tcbox[on line,colback=white,colframe=brand,boxrule=0.7pt,
  arc=3pt,left=4pt,right=4pt,top=0pt,bottom=0pt,boxsep=1pt]{\scriptsize\bfseries\color{brand}#1}}

% ---- Multiple-choice options ----
\tcbset{optbase/.style={enhanced,arc=6pt,boxrule=0.7pt,colframe=linegray,
  left=6pt,right=6pt,top=3pt,bottom=3pt,valign=center}}
\newcommand{\optlabel}[1]{\tikz[baseline=(c.base)]\node[draw=brand,circle,line width=0.8pt,
  inner sep=1.2pt,minimum size=1.35em,font=\footnotesize\bfseries,text=ink](c){\khletter{#1}};}
\newcommand{\opt}[2]{%
  \ifanswers\ifnum#1=\correct\relax
    \begin{tcolorbox}[optbase,colback=rightfill,colframe=brand,boxrule=1.2pt]
      \optlabel{#1}\quad #2\hfill{\color{brand}$\checkmark$}\end{tcolorbox}
  \else
    \begin{tcolorbox}[optbase,colback=white]\optlabel{#1}\quad #2\end{tcolorbox}
  \fi\else
    \begin{tcolorbox}[optbase,colback=white]\optlabel{#1}\quad #2\end{tcolorbox}
  \fi}

% ---- Ruled writing lines (student copy only) ----
\ifanswers
  \tcbset{writelines/.style={}}
\else
  \tcbset{writelines/.style={underlay={\begin{tcbclipinterior}
      \foreach \i in {1,...,30}
        \draw[linegray,line width=0.4pt]
          ($(interior.north west)+(0.5,-\i*0.72+0.05)$) --
          ($(interior.north east)+(-0.5,-\i*0.72+0.05)$);
      \end{tcbclipinterior}}}}
\fi

% ---- Solution text (shown only in teacher copy) ----
\newcommand{\sol}[1]{\gdef\thesol{#1}}

% ---- One exercise = one page ----
\newcounter{prob}
% #1 level  #2 intro  #3 display  #4 question  #5-#8 options  #9 correct (1-4)
\newcommand{\problem}[9]{%
  \stepcounter{prob}\def\correct{#9}%
  \pageheader
  \begin{tcolorbox}[enhanced,colback=mint,colframe=brand,boxrule=0.9pt,arc=8pt,
    left=10pt,right=10pt,top=8pt,bottom=8pt]
    {\bfseries\color{ink}លំហាត់ទី \khnum{\value{prob}}}\enspace\badge{#1}\par\vspace{2pt}
    {\small #2}\par\vspace{6pt}
    \begin{tcolorbox}[enhanced,colback=white,colframe=linegray,boxrule=0.6pt,arc=6pt,
      halign=center,top=6pt,bottom=6pt]
      {\Large\mbox{$\displaystyle #3$}}\par\vspace{4pt}
      {\small\bfseries #4}
    \end{tcolorbox}
  \end{tcolorbox}
  \vspace{6pt}
  {\bfseries ក. គូសរង្វង់ចម្លើយដែលត្រឹមត្រូវ}\par\vspace{4pt}
  \begin{tcbraster}[raster columns=2,raster equal height,raster column skip=12pt,
    raster row skip=7pt]
    \opt{1}{#5}\opt{2}{#6}\opt{3}{#7}\opt{4}{#8}
  \end{tcbraster}
  \vspace{8pt}
  {\bfseries ខ. សរសេរដំណោះស្រាយលម្អិត}\par\vspace{4pt}
  \begin{tcolorbox}[enhanced,height fill,colback=white,colframe=linegray,boxrule=0.6pt,
    arc=8pt,left=14pt,right=14pt,top=10pt,writelines]
    \ifanswers
      {\color{crimson}\bfseries ចម្លើយ ៖ (\khletter{#9})}\par\vspace{2pt}
      {\color{crimson}\thesol}
    \else\mbox{}\fi
  \end{tcolorbox}
  \newpage}

\newcommand{\N}{\,\mathrm{N}}
\newcommand{\kg}{\,\mathrm{kg}}

\begin{document}

% ================= 1 =================
\sol{$F_x=F\cos\alpha=10\times\cos60^\circ=10\times0.5=5\N$}
\problem{ងាយ}
  {កម្លាំង $\vec F$ ធ្វើមុំ $\alpha$ ជាមួយអ័ក្សផ្ដេក $Ox$ ។ គេឲ្យ}
  {F=10\N\,,\qquad \alpha=60^\circ}
  {គណនាកម្លាំងបំបែកផ្ដេក $F_x$}
  {$5\N$}{$8.66\N$}{$10\N$}{$0\N$}{1}

% ================= 2 =================
\sol{$F_y=F\sin\alpha=10\times\sin60^\circ=10\times0.866\approx8.66\N$}
\problem{ងាយ}
  {កម្លាំង $\vec F$ ធ្វើមុំ $\alpha$ ជាមួយអ័ក្សផ្ដេក $Ox$ ។ គេឲ្យ}
  {F=10\N\,,\qquad \alpha=60^\circ}
  {គណនាកម្លាំងបំបែកបញ្ឈរ $F_y$}
  {$5\N$}{$10\N$}{$8.66\N$}{$17.3\N$}{3}

% ================= 3 =================
\sol{$F_y=F\sin\alpha=20\times\sin30^\circ=20\times0.5=10\N$}
\problem{ងាយ}
  {កម្លាំង $\vec F$ ធ្វើមុំ $\alpha$ ជាមួយអ័ក្សផ្ដេក $Ox$ ។ គេឲ្យ}
  {F=20\N\,,\qquad \alpha=30^\circ}
  {គណនាកម្លាំងបំបែកបញ្ឈរ $F_y$}
  {$17.3\N$}{$10\N$}{$20\N$}{$40\N$}{2}

% ================= 4 =================
\sol{$F_x=F\cos45^\circ=20\times0.707\approx14.1\N$}
\problem{ងាយ}
  {កម្លាំង $\vec F$ ធ្វើមុំ $\alpha$ ជាមួយអ័ក្សផ្ដេក $Ox$ ។ គេឲ្យ}
  {F=20\N\,,\qquad \alpha=45^\circ}
  {គណនាកម្លាំងបំបែកផ្ដេក $F_x$}
  {$10\N$}{$20\N$}{$28.3\N$}{$14.1\N$}{4}

% ================= 5 =================
\sol{$F_x=F\cos37^\circ=50\times0.8=40\N$}
\problem{ងាយ}
  {កម្លាំង $\vec F$ ធ្វើមុំ $\alpha$ ជាមួយអ័ក្សផ្ដេក ។ (យក $\sin37^\circ=0.6$ ,
   $\cos37^\circ=0.8$) ។ គេឲ្យ}
  {F=50\N\,,\qquad \alpha=37^\circ}
  {គណនាកម្លាំងបំបែកផ្ដេក $F_x$}
  {$40\N$}{$30\N$}{$50\N$}{$25\N$}{1}

% ================= 6 =================
\sol{$F_y=F\sin37^\circ=50\times0.6=30\N$}
\problem{ងាយ}
  {កម្លាំង $\vec F$ ធ្វើមុំ $\alpha$ ជាមួយអ័ក្សផ្ដេក ។ (យក $\sin37^\circ=0.6$ ,
   $\cos37^\circ=0.8$) ។ គេឲ្យ}
  {F=50\N\,,\qquad \alpha=37^\circ}
  {គណនាកម្លាំងបំបែកបញ្ឈរ $F_y$}
  {$40\N$}{$20\N$}{$30\N$}{$50\N$}{3}

% ================= 7 =================
\sol{$F=\sqrt{F_x^2+F_y^2}=\sqrt{3^2+4^2}=\sqrt{25}=5\N$}
\problem{ងាយ}
  {កម្លាំង $\vec F$ ត្រូវបានបំបែកជាកម្លាំងពីរកែងគ្នា ។ គេឲ្យ}
  {F_x=3\N\,,\qquad F_y=4\N}
  {គណនាអាំងតង់ស៊ីតេនៃកម្លាំងផ្គុំ $F$}
  {$7\N$}{$5\N$}{$1\N$}{$12\N$}{2}

% ================= 8 =================
\sol{$F=\sqrt{6^2+8^2}=\sqrt{100}=10\N$}
\problem{ងាយ}
  {កម្លាំង $\vec F$ ត្រូវបានបំបែកជាកម្លាំងពីរកែងគ្នា ។ គេឲ្យ}
  {F_x=6\N\,,\qquad F_y=8\N}
  {គណនាអាំងតង់ស៊ីតេនៃកម្លាំងផ្គុំ $F$}
  {$14\N$}{$2\N$}{$48\N$}{$10\N$}{4}

% ================= 9 =================
\sol{$F=\sqrt{12^2+5^2}=\sqrt{144+25}=\sqrt{169}=13\N$}
\problem{មធ្យម}
  {កម្លាំង $\vec F$ ត្រូវបានបំបែកជាកម្លាំងពីរកែងគ្នា ។ គេឲ្យ}
  {F_x=12\N\,,\qquad F_y=5\N}
  {គណនាអាំងតង់ស៊ីតេនៃកម្លាំងផ្គុំ $F$}
  {$17\N$}{$7\N$}{$13\N$}{$60\N$}{3}

% ================= 10 =================
\sol{$\tan\alpha=\dfrac{F_y}{F_x}=\dfrac{10}{10}=1\;\Rightarrow\;\alpha=45^\circ$}
\problem{មធ្យម}
  {កម្លាំង $\vec F$ មានកម្លាំងបំបែកកែងគ្នាពីរ ។ គេឲ្យ}
  {F_x=10\N\,,\qquad F_y=10\N}
  {រកមុំ $\alpha$ ដែល $\vec F$ ធ្វើជាមួយអ័ក្សផ្ដេក $Ox$}
  {$45^\circ$}{$30^\circ$}{$60^\circ$}{$90^\circ$}{1}

% ================= 11 =================
\sol{$\tan\alpha=\dfrac{F_y}{F_x}=\dfrac{10}{10\sqrt3}=\dfrac{1}{\sqrt3}\;\Rightarrow\;\alpha=30^\circ$}
\problem{មធ្យម}
  {កម្លាំង $\vec F$ មានកម្លាំងបំបែកកែងគ្នាពីរ ។ គេឲ្យ}
  {F_x=10\sqrt3\N\,,\qquad F_y=10\N}
  {រកមុំ $\alpha$ ដែល $\vec F$ ធ្វើជាមួយអ័ក្សផ្ដេក $Ox$}
  {$60^\circ$}{$30^\circ$}{$45^\circ$}{$90^\circ$}{2}

% ================= 12 =================
\sol{$\cos\alpha=\dfrac{F_x}{F}=\dfrac{20}{40}=0.5\;\Rightarrow\;\alpha=60^\circ$}
\problem{មធ្យម}
  {កម្លាំង $\vec F$ ធ្វើមុំ $\alpha$ ជាមួយអ័ក្សផ្ដេក $Ox$ ។ គេឲ្យ}
  {F=40\N\,,\qquad F_x=20\N}
  {រកមុំ $\alpha$}
  {$30^\circ$}{$45^\circ$}{$90^\circ$}{$60^\circ$}{4}

% ================= 13 =================
\sol{$F_x=F\cos\alpha\;\Rightarrow\;F=\dfrac{F_x}{\cos60^\circ}=\dfrac{15}{0.5}=30\N$}
\problem{មធ្យម}
  {កម្លាំង $\vec F$ ធ្វើមុំ $\alpha$ ជាមួយអ័ក្សផ្ដេក $Ox$ ។ គេឲ្យ}
  {F_x=15\N\,,\qquad \alpha=60^\circ}
  {រកអាំងតង់ស៊ីតេនៃកម្លាំង $F$}
  {$7.5\N$}{$15\N$}{$30\N$}{$26\N$}{3}

% ================= 14 =================
\sol{មុំជាមួយអ័ក្សផ្ដេក $\alpha=90^\circ-60^\circ=30^\circ$ ។\par
     $F_x=F\cos30^\circ=30\times0.866\approx26\N$}
\problem{មធ្យម}
  {កម្លាំង $\vec F$ ធ្វើមុំ $\beta$ ជាមួយអ័ក្ស\textbf{បញ្ឈរ} $Oy$ ។ គេឲ្យ}
  {F=30\N\,,\qquad \beta=60^\circ}
  {គណនាកម្លាំងបំបែកផ្ដេក $F_x$}
  {$26\N$}{$15\N$}{$30\N$}{$52\N$}{1}

% ================= 15 =================
\sol{កម្លាំងទាញទៅមុខ គឺកម្លាំងបំបែកផ្ដេក ៖\par
     $F_x=F\cos60^\circ=40\times0.5=20\N$}
\problem{មធ្យម}
  {សិស្សម្នាក់ទាញរទេះដោយខ្សែ ដែលធ្វើមុំ $\alpha$ ជាមួយផ្ទៃដីផ្ដេក ។ គេឲ្យ}
  {F=40\N\,,\qquad \alpha=60^\circ}
  {រកកម្លាំងដែលទាញរទេះទៅមុខ (កម្លាំងបំបែកផ្ដេក)}
  {$34.6\N$}{$40\N$}{$80\N$}{$20\N$}{4}

% ================= 16 =================
\sol{$P=mg=5\times10=50\N$\par
     $P_{\parallel}=P\sin30^\circ=50\times0.5=25\N$}
\problem{មធ្យម}
  {ប្រអប់ម៉ាស $m$ ស្ថិតលើប្លង់ទេរ ដែលមានមុំទេរ $\alpha$ ។ ($g=10\,\mathrm{N/kg}$) ។ គេឲ្យ}
  {m=5\kg\,,\qquad \alpha=30^\circ}
  {គណនាកម្លាំងបំបែកនៃទម្ងន់ ស្របនឹងប្លង់ទេរ $P_{\parallel}$}
  {$43.3\N$}{$25\N$}{$50\N$}{$5\N$}{2}

% ================= 17 =================
\sol{$P=mg=5\times10=50\N$\par
     $P_{\perp}=P\cos30^\circ=50\times0.866\approx43.3\N$}
\problem{មធ្យម}
  {ប្រអប់ម៉ាស $m$ ស្ថិតលើប្លង់ទេរ ដែលមានមុំទេរ $\alpha$ ។ ($g=10\,\mathrm{N/kg}$) ។ គេឲ្យ}
  {m=5\kg\,,\qquad \alpha=30^\circ}
  {គណនាកម្លាំងបំបែកនៃទម្ងន់ កែងនឹងប្លង់ទេរ $P_{\perp}$}
  {$25\N$}{$50\N$}{$43.3\N$}{$86.6\N$}{3}

% ================= 18 =================
\sol{$P=mg=2\times10=20\N$\par
     $P_{\parallel}=P\sin37^\circ=20\times0.6=12\N$}
\problem{មធ្យម}
  {ប្រអប់ម៉ាស $m$ ស្ថិតលើប្លង់ទេរ ដែលមានមុំទេរ $\alpha$ ។ ($g=10\,\mathrm{N/kg}$ ,
   $\sin37^\circ=0.6$ , $\cos37^\circ=0.8$) ។ គេឲ្យ}
  {m=2\kg\,,\qquad \alpha=37^\circ}
  {គណនាកម្លាំងបំបែកនៃទម្ងន់ ស្របនឹងប្លង់ទេរ $P_{\parallel}$}
  {$12\N$}{$16\N$}{$20\N$}{$6\N$}{1}

% ================= 19 =================
\sol{$P_{\parallel}$ និង $P_{\perp}$ កែងគ្នា ដូច្នេះ\par
     $P=\sqrt{P_{\parallel}^2+P_{\perp}^2}=\sqrt{30^2+40^2}=\sqrt{2500}=50\N$}
\problem{ពិបាក}
  {ទម្ងន់ $\vec P$ របស់វត្ថុមួយលើប្លង់ទេរ ត្រូវបានបំបែកជាកម្លាំងពីរ ។ គេឲ្យ}
  {P_{\parallel}=30\N\,,\qquad P_{\perp}=40\N}
  {រកទម្ងន់ $P$ របស់វត្ថុ}
  {$70\N$}{$10\N$}{$35\N$}{$50\N$}{4}

% ================= 20 =================
\sol{ដើម្បីទប់ប្រអប់ $F$ ត្រូវស្មើនឹង $P_{\parallel}$ ៖\par
     $F=P_{\parallel}=mg\sin30^\circ=4\times10\times0.5=20\N$}
\problem{ពិបាក}
  {ប្រអប់ម៉ាស $m$ ស្ថិតលើប្លង់ទេររលោង (គ្មានកកិត) ដែលមានមុំទេរ $\alpha$ ។
   ដើម្បីទប់ប្រអប់កុំឲ្យរអិល គេអនុវត្តកម្លាំង $\vec F$ ស្របនឹងប្លង់ ។
   ($g=10\,\mathrm{N/kg}$) ។ គេឲ្យ}
  {m=4\kg\,,\qquad \alpha=30^\circ}
  {រកអាំងតង់ស៊ីតេនៃកម្លាំង $F$}
  {$40\N$}{$20\N$}{$34.6\N$}{$10\N$}{2}

\end{document}
"
%  Requires the Khmer font "Noto Sans Khmer".
% =====================================================================
\documentclass[11pt,a4paper]{article}

\usepackage{amsmath,amssymb}
\usepackage{fontspec}
\usepackage[a4paper,top=1.4cm,bottom=1.9cm,left=1.8cm,right=1.8cm,
  footskip=0.9cm]{geometry}
\usepackage{xcolor}
\usepackage{tikz}
\usetikzlibrary{calc}
\usepackage[most]{tcolorbox}
\usepackage{fancyhdr}

% ---- Fonts: Latin = Latin Modern Sans, Khmer = Noto Sans Khmer ----
\setmainfont{Latin Modern Sans}
\newfontfamily\khmerfont[Script=Khmer]{Noto Sans Khmer}
\usepackage[Latin,Khmer]{ucharclasses}
\setTransitionsFor{Khmer}{\khmerfont}{\rmfamily}
\linespread{1.25}
\setlength{\parindent}{0pt}

% ---- Colours ----
\definecolor{brand}{HTML}{1F6B5E}
\definecolor{mint}{HTML}{EEF6F3}
\definecolor{ink}{HTML}{1E2A2A}
\definecolor{inkgray}{HTML}{5D6B6B}
\definecolor{linegray}{HTML}{C9D6D3}
\definecolor{rightfill}{HTML}{DDF1E6}
\definecolor{crimson}{HTML}{B3263A}

% ---- Student / teacher mode ----
\newif\ifanswers
\ifdefined\showanswers \answerstrue \fi

\def\W{17.4} % text width in cm

% ---- Khmer numerals ----
\newcommand{\khnum}[1]{\ifcase#1\relax ០\or ១\or ២\or ៣\or ៤\or ៥\or ៦\or ៧\or
  ៨\or ៩\or ១០\or ១១\or ១២\or ១៣\or ១៤\or ១៥\or ១៦\or ១៧\or ១៨\or ១៩\or ២០\fi}
\newcommand{\khletter}[1]{\ifcase#1\relax\or ក\or ខ\or គ\or ឃ\fi}

% ---- Footer ----
\pagestyle{fancy}
\fancyhf{}
\renewcommand{\headrulewidth}{0pt}
\renewcommand{\footrulewidth}{0.5pt}
\renewcommand{\footrule}{{\color{linegray}\hrule height\footrulewidth width\headwidth}\vspace{3pt}}
\fancyfoot[L]{\scriptsize\color{inkgray}Science Master Academic · រូបវិទ្យា ថ្នាក់ទី០៩ ·
  ទំព័រ \khnum{\value{page}} / ២០}
\fancyfoot[R]{\scriptsize\color{inkgray}ទំនាក់ទំនង 088 909 3291 · Telegram @POVKIMYAT}

% ---- Page header ----
\newcommand{\pageheader}{%
\begin{tikzpicture}
  \useasboundingbox (0,0.1) rectangle (\W,-2.45);
  % logo
  \begin{scope}[shift={(0.9,-1.05)}]
    \fill[mint] (0,0) circle (0.88);
    \draw[brand,line width=1.1pt] (0,0) circle (0.88);
    \foreach \x/\y in {-0.19/0.36,0.19/0.36,-0.19/-0.02,0.19/-0.02}
      \draw[brand,line width=1.1pt,rounded corners=1.5pt]
        (\x-0.15,\y-0.15) rectangle (\x+0.15,\y+0.15);
    \node[font=\fontsize{5}{5}\selectfont\bfseries,text=brand] at (0,-0.5) {Science Master};
  \end{scope}
  % titles
  \node[anchor=west,font=\small\bfseries,text=brand] at (2.05,-0.3) {SCIENCE MASTER ACADEMIC};
  \node[anchor=west,font=\LARGE\bfseries,text=ink] at (2.05,-1.0)
    {សន្លឹកកិច្ចការ\ifanswers{\large\color{crimson}\ (ចម្លើយសម្រាប់គ្រូ)}\fi};
  \node[anchor=west,font=\small\bfseries,text=ink] at (2.05,-1.75)
    {រូបវិទ្យា ថ្នាក់ទី០៩ · មេរៀនទី០១ ៖ ការបំបែកកម្លាំង};
  % score box
  \draw[linegray,rounded corners=6pt,line width=0.9pt] (\W-2.6,-0.15) rectangle (\W,-1.75);
  \node[font=\footnotesize,text=inkgray] at (\W-1.3,-0.45) {ពិន្ទុ};
  \node[font=\large,text=ink] at (\W-1.3,-1.2) {\rule{1cm}{0.5pt}\, / 10};
  % rule
  \draw[brand,line width=2pt] (0,-2.35) -- (\W,-2.35);
\end{tikzpicture}\par\vspace{8pt}
{\small ឈ្មោះសិស្ស ៖ \makebox[6.2cm]{\color{inkgray}\dotfill}\hfill
 កាលបរិច្ឆេទ ៖ \makebox[4.2cm]{\color{inkgray}\dotfill}}\par\vspace{8pt}}

% ---- Difficulty badge ----
\newcommand{\badge}[1]{\tcbox[on line,colback=white,colframe=brand,boxrule=0.7pt,
  arc=3pt,left=4pt,right=4pt,top=0pt,bottom=0pt,boxsep=1pt]{\scriptsize\bfseries\color{brand}#1}}

% ---- Multiple-choice options ----
\tcbset{optbase/.style={enhanced,arc=6pt,boxrule=0.7pt,colframe=linegray,
  left=6pt,right=6pt,top=3pt,bottom=3pt,valign=center}}
\newcommand{\optlabel}[1]{\tikz[baseline=(c.base)]\node[draw=brand,circle,line width=0.8pt,
  inner sep=1.2pt,minimum size=1.35em,font=\footnotesize\bfseries,text=ink](c){\khletter{#1}};}
\newcommand{\opt}[2]{%
  \ifanswers\ifnum#1=\correct\relax
    \begin{tcolorbox}[optbase,colback=rightfill,colframe=brand,boxrule=1.2pt]
      \optlabel{#1}\quad #2\hfill{\color{brand}$\checkmark$}\end{tcolorbox}
  \else
    \begin{tcolorbox}[optbase,colback=white]\optlabel{#1}\quad #2\end{tcolorbox}
  \fi\else
    \begin{tcolorbox}[optbase,colback=white]\optlabel{#1}\quad #2\end{tcolorbox}
  \fi}

% ---- Ruled writing lines (student copy only) ----
\ifanswers
  \tcbset{writelines/.style={}}
\else
  \tcbset{writelines/.style={underlay={\begin{tcbclipinterior}
      \foreach \i in {1,...,30}
        \draw[linegray,line width=0.4pt]
          ($(interior.north west)+(0.5,-\i*0.72+0.05)$) --
          ($(interior.north east)+(-0.5,-\i*0.72+0.05)$);
      \end{tcbclipinterior}}}}
\fi

% ---- Solution text (shown only in teacher copy) ----
\newcommand{\sol}[1]{\gdef\thesol{#1}}

% ---- One exercise = one page ----
\newcounter{prob}
% #1 level  #2 intro  #3 display  #4 question  #5-#8 options  #9 correct (1-4)
\newcommand{\problem}[9]{%
  \stepcounter{prob}\def\correct{#9}%
  \pageheader
  \begin{tcolorbox}[enhanced,colback=mint,colframe=brand,boxrule=0.9pt,arc=8pt,
    left=10pt,right=10pt,top=8pt,bottom=8pt]
    {\bfseries\color{ink}លំហាត់ទី \khnum{\value{prob}}}\enspace\badge{#1}\par\vspace{2pt}
    {\small #2}\par\vspace{6pt}
    \begin{tcolorbox}[enhanced,colback=white,colframe=linegray,boxrule=0.6pt,arc=6pt,
      halign=center,top=6pt,bottom=6pt]
      {\Large\mbox{$\displaystyle #3$}}\par\vspace{4pt}
      {\small\bfseries #4}
    \end{tcolorbox}
  \end{tcolorbox}
  \vspace{6pt}
  {\bfseries ក. គូសរង្វង់ចម្លើយដែលត្រឹមត្រូវ}\par\vspace{4pt}
  \begin{tcbraster}[raster columns=2,raster equal height,raster column skip=12pt,
    raster row skip=7pt]
    \opt{1}{#5}\opt{2}{#6}\opt{3}{#7}\opt{4}{#8}
  \end{tcbraster}
  \vspace{8pt}
  {\bfseries ខ. សរសេរដំណោះស្រាយលម្អិត}\par\vspace{4pt}
  \begin{tcolorbox}[enhanced,height fill,colback=white,colframe=linegray,boxrule=0.6pt,
    arc=8pt,left=14pt,right=14pt,top=10pt,writelines]
    \ifanswers
      {\color{crimson}\bfseries ចម្លើយ ៖ (\khletter{#9})}\par\vspace{2pt}
      {\color{crimson}\thesol}
    \else\mbox{}\fi
  \end{tcolorbox}
  \newpage}

\newcommand{\N}{\,\mathrm{N}}
\newcommand{\kg}{\,\mathrm{kg}}

\begin{document}

% ================= 1 =================
\sol{$F_x=F\cos\alpha=10\times\cos60^\circ=10\times0.5=5\N$}
\problem{ងាយ}
  {កម្លាំង $\vec F$ ធ្វើមុំ $\alpha$ ជាមួយអ័ក្សផ្ដេក $Ox$ ។ គេឲ្យ}
  {F=10\N\,,\qquad \alpha=60^\circ}
  {គណនាកម្លាំងបំបែកផ្ដេក $F_x$}
  {$5\N$}{$8.66\N$}{$10\N$}{$0\N$}{1}

% ================= 2 =================
\sol{$F_y=F\sin\alpha=10\times\sin60^\circ=10\times0.866\approx8.66\N$}
\problem{ងាយ}
  {កម្លាំង $\vec F$ ធ្វើមុំ $\alpha$ ជាមួយអ័ក្សផ្ដេក $Ox$ ។ គេឲ្យ}
  {F=10\N\,,\qquad \alpha=60^\circ}
  {គណនាកម្លាំងបំបែកបញ្ឈរ $F_y$}
  {$5\N$}{$10\N$}{$8.66\N$}{$17.3\N$}{3}

% ================= 3 =================
\sol{$F_y=F\sin\alpha=20\times\sin30^\circ=20\times0.5=10\N$}
\problem{ងាយ}
  {កម្លាំង $\vec F$ ធ្វើមុំ $\alpha$ ជាមួយអ័ក្សផ្ដេក $Ox$ ។ គេឲ្យ}
  {F=20\N\,,\qquad \alpha=30^\circ}
  {គណនាកម្លាំងបំបែកបញ្ឈរ $F_y$}
  {$17.3\N$}{$10\N$}{$20\N$}{$40\N$}{2}

% ================= 4 =================
\sol{$F_x=F\cos45^\circ=20\times0.707\approx14.1\N$}
\problem{ងាយ}
  {កម្លាំង $\vec F$ ធ្វើមុំ $\alpha$ ជាមួយអ័ក្សផ្ដេក $Ox$ ។ គេឲ្យ}
  {F=20\N\,,\qquad \alpha=45^\circ}
  {គណនាកម្លាំងបំបែកផ្ដេក $F_x$}
  {$10\N$}{$20\N$}{$28.3\N$}{$14.1\N$}{4}

% ================= 5 =================
\sol{$F_x=F\cos37^\circ=50\times0.8=40\N$}
\problem{ងាយ}
  {កម្លាំង $\vec F$ ធ្វើមុំ $\alpha$ ជាមួយអ័ក្សផ្ដេក ។ (យក $\sin37^\circ=0.6$ ,
   $\cos37^\circ=0.8$) ។ គេឲ្យ}
  {F=50\N\,,\qquad \alpha=37^\circ}
  {គណនាកម្លាំងបំបែកផ្ដេក $F_x$}
  {$40\N$}{$30\N$}{$50\N$}{$25\N$}{1}

% ================= 6 =================
\sol{$F_y=F\sin37^\circ=50\times0.6=30\N$}
\problem{ងាយ}
  {កម្លាំង $\vec F$ ធ្វើមុំ $\alpha$ ជាមួយអ័ក្សផ្ដេក ។ (យក $\sin37^\circ=0.6$ ,
   $\cos37^\circ=0.8$) ។ គេឲ្យ}
  {F=50\N\,,\qquad \alpha=37^\circ}
  {គណនាកម្លាំងបំបែកបញ្ឈរ $F_y$}
  {$40\N$}{$20\N$}{$30\N$}{$50\N$}{3}

% ================= 7 =================
\sol{$F=\sqrt{F_x^2+F_y^2}=\sqrt{3^2+4^2}=\sqrt{25}=5\N$}
\problem{ងាយ}
  {កម្លាំង $\vec F$ ត្រូវបានបំបែកជាកម្លាំងពីរកែងគ្នា ។ គេឲ្យ}
  {F_x=3\N\,,\qquad F_y=4\N}
  {គណនាអាំងតង់ស៊ីតេនៃកម្លាំងផ្គុំ $F$}
  {$7\N$}{$5\N$}{$1\N$}{$12\N$}{2}

% ================= 8 =================
\sol{$F=\sqrt{6^2+8^2}=\sqrt{100}=10\N$}
\problem{ងាយ}
  {កម្លាំង $\vec F$ ត្រូវបានបំបែកជាកម្លាំងពីរកែងគ្នា ។ គេឲ្យ}
  {F_x=6\N\,,\qquad F_y=8\N}
  {គណនាអាំងតង់ស៊ីតេនៃកម្លាំងផ្គុំ $F$}
  {$14\N$}{$2\N$}{$48\N$}{$10\N$}{4}

% ================= 9 =================
\sol{$F=\sqrt{12^2+5^2}=\sqrt{144+25}=\sqrt{169}=13\N$}
\problem{មធ្យម}
  {កម្លាំង $\vec F$ ត្រូវបានបំបែកជាកម្លាំងពីរកែងគ្នា ។ គេឲ្យ}
  {F_x=12\N\,,\qquad F_y=5\N}
  {គណនាអាំងតង់ស៊ីតេនៃកម្លាំងផ្គុំ $F$}
  {$17\N$}{$7\N$}{$13\N$}{$60\N$}{3}

% ================= 10 =================
\sol{$\tan\alpha=\dfrac{F_y}{F_x}=\dfrac{10}{10}=1\;\Rightarrow\;\alpha=45^\circ$}
\problem{មធ្យម}
  {កម្លាំង $\vec F$ មានកម្លាំងបំបែកកែងគ្នាពីរ ។ គេឲ្យ}
  {F_x=10\N\,,\qquad F_y=10\N}
  {រកមុំ $\alpha$ ដែល $\vec F$ ធ្វើជាមួយអ័ក្សផ្ដេក $Ox$}
  {$45^\circ$}{$30^\circ$}{$60^\circ$}{$90^\circ$}{1}

% ================= 11 =================
\sol{$\tan\alpha=\dfrac{F_y}{F_x}=\dfrac{10}{10\sqrt3}=\dfrac{1}{\sqrt3}\;\Rightarrow\;\alpha=30^\circ$}
\problem{មធ្យម}
  {កម្លាំង $\vec F$ មានកម្លាំងបំបែកកែងគ្នាពីរ ។ គេឲ្យ}
  {F_x=10\sqrt3\N\,,\qquad F_y=10\N}
  {រកមុំ $\alpha$ ដែល $\vec F$ ធ្វើជាមួយអ័ក្សផ្ដេក $Ox$}
  {$60^\circ$}{$30^\circ$}{$45^\circ$}{$90^\circ$}{2}

% ================= 12 =================
\sol{$\cos\alpha=\dfrac{F_x}{F}=\dfrac{20}{40}=0.5\;\Rightarrow\;\alpha=60^\circ$}
\problem{មធ្យម}
  {កម្លាំង $\vec F$ ធ្វើមុំ $\alpha$ ជាមួយអ័ក្សផ្ដេក $Ox$ ។ គេឲ្យ}
  {F=40\N\,,\qquad F_x=20\N}
  {រកមុំ $\alpha$}
  {$30^\circ$}{$45^\circ$}{$90^\circ$}{$60^\circ$}{4}

% ================= 13 =================
\sol{$F_x=F\cos\alpha\;\Rightarrow\;F=\dfrac{F_x}{\cos60^\circ}=\dfrac{15}{0.5}=30\N$}
\problem{មធ្យម}
  {កម្លាំង $\vec F$ ធ្វើមុំ $\alpha$ ជាមួយអ័ក្សផ្ដេក $Ox$ ។ គេឲ្យ}
  {F_x=15\N\,,\qquad \alpha=60^\circ}
  {រកអាំងតង់ស៊ីតេនៃកម្លាំង $F$}
  {$7.5\N$}{$15\N$}{$30\N$}{$26\N$}{3}

% ================= 14 =================
\sol{មុំជាមួយអ័ក្សផ្ដេក $\alpha=90^\circ-60^\circ=30^\circ$ ។\par
     $F_x=F\cos30^\circ=30\times0.866\approx26\N$}
\problem{មធ្យម}
  {កម្លាំង $\vec F$ ធ្វើមុំ $\beta$ ជាមួយអ័ក្ស\textbf{បញ្ឈរ} $Oy$ ។ គេឲ្យ}
  {F=30\N\,,\qquad \beta=60^\circ}
  {គណនាកម្លាំងបំបែកផ្ដេក $F_x$}
  {$26\N$}{$15\N$}{$30\N$}{$52\N$}{1}

% ================= 15 =================
\sol{កម្លាំងទាញទៅមុខ គឺកម្លាំងបំបែកផ្ដេក ៖\par
     $F_x=F\cos60^\circ=40\times0.5=20\N$}
\problem{មធ្យម}
  {សិស្សម្នាក់ទាញរទេះដោយខ្សែ ដែលធ្វើមុំ $\alpha$ ជាមួយផ្ទៃដីផ្ដេក ។ គេឲ្យ}
  {F=40\N\,,\qquad \alpha=60^\circ}
  {រកកម្លាំងដែលទាញរទេះទៅមុខ (កម្លាំងបំបែកផ្ដេក)}
  {$34.6\N$}{$40\N$}{$80\N$}{$20\N$}{4}

% ================= 16 =================
\sol{$P=mg=5\times10=50\N$\par
     $P_{\parallel}=P\sin30^\circ=50\times0.5=25\N$}
\problem{មធ្យម}
  {ប្រអប់ម៉ាស $m$ ស្ថិតលើប្លង់ទេរ ដែលមានមុំទេរ $\alpha$ ។ ($g=10\,\mathrm{N/kg}$) ។ គេឲ្យ}
  {m=5\kg\,,\qquad \alpha=30^\circ}
  {គណនាកម្លាំងបំបែកនៃទម្ងន់ ស្របនឹងប្លង់ទេរ $P_{\parallel}$}
  {$43.3\N$}{$25\N$}{$50\N$}{$5\N$}{2}

% ================= 17 =================
\sol{$P=mg=5\times10=50\N$\par
     $P_{\perp}=P\cos30^\circ=50\times0.866\approx43.3\N$}
\problem{មធ្យម}
  {ប្រអប់ម៉ាស $m$ ស្ថិតលើប្លង់ទេរ ដែលមានមុំទេរ $\alpha$ ។ ($g=10\,\mathrm{N/kg}$) ។ គេឲ្យ}
  {m=5\kg\,,\qquad \alpha=30^\circ}
  {គណនាកម្លាំងបំបែកនៃទម្ងន់ កែងនឹងប្លង់ទេរ $P_{\perp}$}
  {$25\N$}{$50\N$}{$43.3\N$}{$86.6\N$}{3}

% ================= 18 =================
\sol{$P=mg=2\times10=20\N$\par
     $P_{\parallel}=P\sin37^\circ=20\times0.6=12\N$}
\problem{មធ្យម}
  {ប្រអប់ម៉ាស $m$ ស្ថិតលើប្លង់ទេរ ដែលមានមុំទេរ $\alpha$ ។ ($g=10\,\mathrm{N/kg}$ ,
   $\sin37^\circ=0.6$ , $\cos37^\circ=0.8$) ។ គេឲ្យ}
  {m=2\kg\,,\qquad \alpha=37^\circ}
  {គណនាកម្លាំងបំបែកនៃទម្ងន់ ស្របនឹងប្លង់ទេរ $P_{\parallel}$}
  {$12\N$}{$16\N$}{$20\N$}{$6\N$}{1}

% ================= 19 =================
\sol{$P_{\parallel}$ និង $P_{\perp}$ កែងគ្នា ដូច្នេះ\par
     $P=\sqrt{P_{\parallel}^2+P_{\perp}^2}=\sqrt{30^2+40^2}=\sqrt{2500}=50\N$}
\problem{ពិបាក}
  {ទម្ងន់ $\vec P$ របស់វត្ថុមួយលើប្លង់ទេរ ត្រូវបានបំបែកជាកម្លាំងពីរ ។ គេឲ្យ}
  {P_{\parallel}=30\N\,,\qquad P_{\perp}=40\N}
  {រកទម្ងន់ $P$ របស់វត្ថុ}
  {$70\N$}{$10\N$}{$35\N$}{$50\N$}{4}

% ================= 20 =================
\sol{ដើម្បីទប់ប្រអប់ $F$ ត្រូវស្មើនឹង $P_{\parallel}$ ៖\par
     $F=P_{\parallel}=mg\sin30^\circ=4\times10\times0.5=20\N$}
\problem{ពិបាក}
  {ប្រអប់ម៉ាស $m$ ស្ថិតលើប្លង់ទេររលោង (គ្មានកកិត) ដែលមានមុំទេរ $\alpha$ ។
   ដើម្បីទប់ប្រអប់កុំឲ្យរអិល គេអនុវត្តកម្លាំង $\vec F$ ស្របនឹងប្លង់ ។
   ($g=10\,\mathrm{N/kg}$) ។ គេឲ្យ}
  {m=4\kg\,,\qquad \alpha=30^\circ}
  {រកអាំងតង់ស៊ីតេនៃកម្លាំង $F$}
  {$40\N$}{$20\N$}{$34.6\N$}{$10\N$}{2}

\end{document}
"
%  Requires the Khmer font "Noto Sans Khmer".
% =====================================================================
\documentclass[11pt,a4paper]{article}

\usepackage{amsmath,amssymb}
\usepackage{fontspec}
\usepackage[a4paper,top=1.4cm,bottom=1.9cm,left=1.8cm,right=1.8cm,
  footskip=0.9cm]{geometry}
\usepackage{xcolor}
\usepackage{tikz}
\usetikzlibrary{calc}
\usepackage[most]{tcolorbox}
\usepackage{fancyhdr}

% ---- Fonts: Latin = Latin Modern Sans, Khmer = Noto Sans Khmer ----
\setmainfont{Latin Modern Sans}
\newfontfamily\khmerfont[Script=Khmer]{Noto Sans Khmer}
\usepackage[Latin,Khmer]{ucharclasses}
\setTransitionsFor{Khmer}{\khmerfont}{\rmfamily}
\linespread{1.25}
\setlength{\parindent}{0pt}

% ---- Colours ----
\definecolor{brand}{HTML}{1F6B5E}
\definecolor{mint}{HTML}{EEF6F3}
\definecolor{ink}{HTML}{1E2A2A}
\definecolor{inkgray}{HTML}{5D6B6B}
\definecolor{linegray}{HTML}{C9D6D3}
\definecolor{rightfill}{HTML}{DDF1E6}
\definecolor{crimson}{HTML}{B3263A}

% ---- Student / teacher mode ----
\newif\ifanswers
\ifdefined\showanswers \answerstrue \fi

\def\W{17.4} % text width in cm

% ---- Khmer numerals ----
\newcommand{\khnum}[1]{\ifcase#1\relax ០\or ១\or ២\or ៣\or ៤\or ៥\or ៦\or ៧\or
  ៨\or ៩\or ១០\or ១១\or ១២\or ១៣\or ១៤\or ១៥\or ១៦\or ១៧\or ១៨\or ១៩\or ២០\fi}
\newcommand{\khletter}[1]{\ifcase#1\relax\or ក\or ខ\or គ\or ឃ\fi}

% ---- Footer ----
\pagestyle{fancy}
\fancyhf{}
\renewcommand{\headrulewidth}{0pt}
\renewcommand{\footrulewidth}{0.5pt}
\renewcommand{\footrule}{{\color{linegray}\hrule height\footrulewidth width\headwidth}\vspace{3pt}}
\fancyfoot[L]{\scriptsize\color{inkgray}Science Master Academic · រូបវិទ្យា ថ្នាក់ទី០៩ ·
  ទំព័រ \khnum{\value{page}} / ២០}
\fancyfoot[R]{\scriptsize\color{inkgray}ទំនាក់ទំនង 088 909 3291 · Telegram @POVKIMYAT}

% ---- Page header ----
\newcommand{\pageheader}{%
\begin{tikzpicture}
  \useasboundingbox (0,0.1) rectangle (\W,-2.45);
  % logo
  \begin{scope}[shift={(0.9,-1.05)}]
    \fill[mint] (0,0) circle (0.88);
    \draw[brand,line width=1.1pt] (0,0) circle (0.88);
    \foreach \x/\y in {-0.19/0.36,0.19/0.36,-0.19/-0.02,0.19/-0.02}
      \draw[brand,line width=1.1pt,rounded corners=1.5pt]
        (\x-0.15,\y-0.15) rectangle (\x+0.15,\y+0.15);
    \node[font=\fontsize{5}{5}\selectfont\bfseries,text=brand] at (0,-0.5) {Science Master};
  \end{scope}
  % titles
  \node[anchor=west,font=\small\bfseries,text=brand] at (2.05,-0.3) {SCIENCE MASTER ACADEMIC};
  \node[anchor=west,font=\LARGE\bfseries,text=ink] at (2.05,-1.0)
    {សន្លឹកកិច្ចការ\ifanswers{\large\color{crimson}\ (ចម្លើយសម្រាប់គ្រូ)}\fi};
  \node[anchor=west,font=\small\bfseries,text=ink] at (2.05,-1.75)
    {រូបវិទ្យា ថ្នាក់ទី០៩ · មេរៀនទី០១ ៖ ការបំបែកកម្លាំង};
  % score box
  \draw[linegray,rounded corners=6pt,line width=0.9pt] (\W-2.6,-0.15) rectangle (\W,-1.75);
  \node[font=\footnotesize,text=inkgray] at (\W-1.3,-0.45) {ពិន្ទុ};
  \node[font=\large,text=ink] at (\W-1.3,-1.2) {\rule{1cm}{0.5pt}\, / 10};
  % rule
  \draw[brand,line width=2pt] (0,-2.35) -- (\W,-2.35);
\end{tikzpicture}\par\vspace{8pt}
{\small ឈ្មោះសិស្ស ៖ \makebox[6.2cm]{\color{inkgray}\dotfill}\hfill
 កាលបរិច្ឆេទ ៖ \makebox[4.2cm]{\color{inkgray}\dotfill}}\par\vspace{8pt}}

% ---- Difficulty badge ----
\newcommand{\badge}[1]{\tcbox[on line,colback=white,colframe=brand,boxrule=0.7pt,
  arc=3pt,left=4pt,right=4pt,top=0pt,bottom=0pt,boxsep=1pt]{\scriptsize\bfseries\color{brand}#1}}

% ---- Multiple-choice options ----
\tcbset{optbase/.style={enhanced,arc=6pt,boxrule=0.7pt,colframe=linegray,
  left=6pt,right=6pt,top=3pt,bottom=3pt,valign=center}}
\newcommand{\optlabel}[1]{\tikz[baseline=(c.base)]\node[draw=brand,circle,line width=0.8pt,
  inner sep=1.2pt,minimum size=1.35em,font=\footnotesize\bfseries,text=ink](c){\khletter{#1}};}
\newcommand{\opt}[2]{%
  \ifanswers\ifnum#1=\correct\relax
    \begin{tcolorbox}[optbase,colback=rightfill,colframe=brand,boxrule=1.2pt]
      \optlabel{#1}\quad #2\hfill{\color{brand}$\checkmark$}\end{tcolorbox}
  \else
    \begin{tcolorbox}[optbase,colback=white]\optlabel{#1}\quad #2\end{tcolorbox}
  \fi\else
    \begin{tcolorbox}[optbase,colback=white]\optlabel{#1}\quad #2\end{tcolorbox}
  \fi}

% ---- Ruled writing lines (student copy only) ----
\ifanswers
  \tcbset{writelines/.style={}}
\else
  \tcbset{writelines/.style={underlay={\begin{tcbclipinterior}
      \foreach \i in {1,...,30}
        \draw[linegray,line width=0.4pt]
          ($(interior.north west)+(0.5,-\i*0.72+0.05)$) --
          ($(interior.north east)+(-0.5,-\i*0.72+0.05)$);
      \end{tcbclipinterior}}}}
\fi

% ---- Solution text (shown only in teacher copy) ----
\newcommand{\sol}[1]{\gdef\thesol{#1}}

% ---- One exercise = one page ----
\newcounter{prob}
% #1 level  #2 intro  #3 display  #4 question  #5-#8 options  #9 correct (1-4)
\newcommand{\problem}[9]{%
  \stepcounter{prob}\def\correct{#9}%
  \pageheader
  \begin{tcolorbox}[enhanced,colback=mint,colframe=brand,boxrule=0.9pt,arc=8pt,
    left=10pt,right=10pt,top=8pt,bottom=8pt]
    {\bfseries\color{ink}លំហាត់ទី \khnum{\value{prob}}}\enspace\badge{#1}\par\vspace{2pt}
    {\small #2}\par\vspace{6pt}
    \begin{tcolorbox}[enhanced,colback=white,colframe=linegray,boxrule=0.6pt,arc=6pt,
      halign=center,top=6pt,bottom=6pt]
      {\Large\mbox{$\displaystyle #3$}}\par\vspace{4pt}
      {\small\bfseries #4}
    \end{tcolorbox}
  \end{tcolorbox}
  \vspace{6pt}
  {\bfseries ក. គូសរង្វង់ចម្លើយដែលត្រឹមត្រូវ}\par\vspace{4pt}
  \begin{tcbraster}[raster columns=2,raster equal height,raster column skip=12pt,
    raster row skip=7pt]
    \opt{1}{#5}\opt{2}{#6}\opt{3}{#7}\opt{4}{#8}
  \end{tcbraster}
  \vspace{8pt}
  {\bfseries ខ. សរសេរដំណោះស្រាយលម្អិត}\par\vspace{4pt}
  \begin{tcolorbox}[enhanced,height fill,colback=white,colframe=linegray,boxrule=0.6pt,
    arc=8pt,left=14pt,right=14pt,top=10pt,writelines]
    \ifanswers
      {\color{crimson}\bfseries ចម្លើយ ៖ (\khletter{#9})}\par\vspace{2pt}
      {\color{crimson}\thesol}
    \else\mbox{}\fi
  \end{tcolorbox}
  \newpage}

\newcommand{\N}{\,\mathrm{N}}
\newcommand{\kg}{\,\mathrm{kg}}

\begin{document}

% ================= 1 =================
\sol{$F_x=F\cos\alpha=10\times\cos60^\circ=10\times0.5=5\N$}
\problem{ងាយ}
  {កម្លាំង $\vec F$ ធ្វើមុំ $\alpha$ ជាមួយអ័ក្សផ្ដេក $Ox$ ។ គេឲ្យ}
  {F=10\N\,,\qquad \alpha=60^\circ}
  {គណនាកម្លាំងបំបែកផ្ដេក $F_x$}
  {$5\N$}{$8.66\N$}{$10\N$}{$0\N$}{1}

% ================= 2 =================
\sol{$F_y=F\sin\alpha=10\times\sin60^\circ=10\times0.866\approx8.66\N$}
\problem{ងាយ}
  {កម្លាំង $\vec F$ ធ្វើមុំ $\alpha$ ជាមួយអ័ក្សផ្ដេក $Ox$ ។ គេឲ្យ}
  {F=10\N\,,\qquad \alpha=60^\circ}
  {គណនាកម្លាំងបំបែកបញ្ឈរ $F_y$}
  {$5\N$}{$10\N$}{$8.66\N$}{$17.3\N$}{3}

% ================= 3 =================
\sol{$F_y=F\sin\alpha=20\times\sin30^\circ=20\times0.5=10\N$}
\problem{ងាយ}
  {កម្លាំង $\vec F$ ធ្វើមុំ $\alpha$ ជាមួយអ័ក្សផ្ដេក $Ox$ ។ គេឲ្យ}
  {F=20\N\,,\qquad \alpha=30^\circ}
  {គណនាកម្លាំងបំបែកបញ្ឈរ $F_y$}
  {$17.3\N$}{$10\N$}{$20\N$}{$40\N$}{2}

% ================= 4 =================
\sol{$F_x=F\cos45^\circ=20\times0.707\approx14.1\N$}
\problem{ងាយ}
  {កម្លាំង $\vec F$ ធ្វើមុំ $\alpha$ ជាមួយអ័ក្សផ្ដេក $Ox$ ។ គេឲ្យ}
  {F=20\N\,,\qquad \alpha=45^\circ}
  {គណនាកម្លាំងបំបែកផ្ដេក $F_x$}
  {$10\N$}{$20\N$}{$28.3\N$}{$14.1\N$}{4}

% ================= 5 =================
\sol{$F_x=F\cos37^\circ=50\times0.8=40\N$}
\problem{ងាយ}
  {កម្លាំង $\vec F$ ធ្វើមុំ $\alpha$ ជាមួយអ័ក្សផ្ដេក ។ (យក $\sin37^\circ=0.6$ ,
   $\cos37^\circ=0.8$) ។ គេឲ្យ}
  {F=50\N\,,\qquad \alpha=37^\circ}
  {គណនាកម្លាំងបំបែកផ្ដេក $F_x$}
  {$40\N$}{$30\N$}{$50\N$}{$25\N$}{1}

% ================= 6 =================
\sol{$F_y=F\sin37^\circ=50\times0.6=30\N$}
\problem{ងាយ}
  {កម្លាំង $\vec F$ ធ្វើមុំ $\alpha$ ជាមួយអ័ក្សផ្ដេក ។ (យក $\sin37^\circ=0.6$ ,
   $\cos37^\circ=0.8$) ។ គេឲ្យ}
  {F=50\N\,,\qquad \alpha=37^\circ}
  {គណនាកម្លាំងបំបែកបញ្ឈរ $F_y$}
  {$40\N$}{$20\N$}{$30\N$}{$50\N$}{3}

% ================= 7 =================
\sol{$F=\sqrt{F_x^2+F_y^2}=\sqrt{3^2+4^2}=\sqrt{25}=5\N$}
\problem{ងាយ}
  {កម្លាំង $\vec F$ ត្រូវបានបំបែកជាកម្លាំងពីរកែងគ្នា ។ គេឲ្យ}
  {F_x=3\N\,,\qquad F_y=4\N}
  {គណនាអាំងតង់ស៊ីតេនៃកម្លាំងផ្គុំ $F$}
  {$7\N$}{$5\N$}{$1\N$}{$12\N$}{2}

% ================= 8 =================
\sol{$F=\sqrt{6^2+8^2}=\sqrt{100}=10\N$}
\problem{ងាយ}
  {កម្លាំង $\vec F$ ត្រូវបានបំបែកជាកម្លាំងពីរកែងគ្នា ។ គេឲ្យ}
  {F_x=6\N\,,\qquad F_y=8\N}
  {គណនាអាំងតង់ស៊ីតេនៃកម្លាំងផ្គុំ $F$}
  {$14\N$}{$2\N$}{$48\N$}{$10\N$}{4}

% ================= 9 =================
\sol{$F=\sqrt{12^2+5^2}=\sqrt{144+25}=\sqrt{169}=13\N$}
\problem{មធ្យម}
  {កម្លាំង $\vec F$ ត្រូវបានបំបែកជាកម្លាំងពីរកែងគ្នា ។ គេឲ្យ}
  {F_x=12\N\,,\qquad F_y=5\N}
  {គណនាអាំងតង់ស៊ីតេនៃកម្លាំងផ្គុំ $F$}
  {$17\N$}{$7\N$}{$13\N$}{$60\N$}{3}

% ================= 10 =================
\sol{$\tan\alpha=\dfrac{F_y}{F_x}=\dfrac{10}{10}=1\;\Rightarrow\;\alpha=45^\circ$}
\problem{មធ្យម}
  {កម្លាំង $\vec F$ មានកម្លាំងបំបែកកែងគ្នាពីរ ។ គេឲ្យ}
  {F_x=10\N\,,\qquad F_y=10\N}
  {រកមុំ $\alpha$ ដែល $\vec F$ ធ្វើជាមួយអ័ក្សផ្ដេក $Ox$}
  {$45^\circ$}{$30^\circ$}{$60^\circ$}{$90^\circ$}{1}

% ================= 11 =================
\sol{$\tan\alpha=\dfrac{F_y}{F_x}=\dfrac{10}{10\sqrt3}=\dfrac{1}{\sqrt3}\;\Rightarrow\;\alpha=30^\circ$}
\problem{មធ្យម}
  {កម្លាំង $\vec F$ មានកម្លាំងបំបែកកែងគ្នាពីរ ។ គេឲ្យ}
  {F_x=10\sqrt3\N\,,\qquad F_y=10\N}
  {រកមុំ $\alpha$ ដែល $\vec F$ ធ្វើជាមួយអ័ក្សផ្ដេក $Ox$}
  {$60^\circ$}{$30^\circ$}{$45^\circ$}{$90^\circ$}{2}

% ================= 12 =================
\sol{$\cos\alpha=\dfrac{F_x}{F}=\dfrac{20}{40}=0.5\;\Rightarrow\;\alpha=60^\circ$}
\problem{មធ្យម}
  {កម្លាំង $\vec F$ ធ្វើមុំ $\alpha$ ជាមួយអ័ក្សផ្ដេក $Ox$ ។ គេឲ្យ}
  {F=40\N\,,\qquad F_x=20\N}
  {រកមុំ $\alpha$}
  {$30^\circ$}{$45^\circ$}{$90^\circ$}{$60^\circ$}{4}

% ================= 13 =================
\sol{$F_x=F\cos\alpha\;\Rightarrow\;F=\dfrac{F_x}{\cos60^\circ}=\dfrac{15}{0.5}=30\N$}
\problem{មធ្យម}
  {កម្លាំង $\vec F$ ធ្វើមុំ $\alpha$ ជាមួយអ័ក្សផ្ដេក $Ox$ ។ គេឲ្យ}
  {F_x=15\N\,,\qquad \alpha=60^\circ}
  {រកអាំងតង់ស៊ីតេនៃកម្លាំង $F$}
  {$7.5\N$}{$15\N$}{$30\N$}{$26\N$}{3}

% ================= 14 =================
\sol{មុំជាមួយអ័ក្សផ្ដេក $\alpha=90^\circ-60^\circ=30^\circ$ ។\par
     $F_x=F\cos30^\circ=30\times0.866\approx26\N$}
\problem{មធ្យម}
  {កម្លាំង $\vec F$ ធ្វើមុំ $\beta$ ជាមួយអ័ក្ស\textbf{បញ្ឈរ} $Oy$ ។ គេឲ្យ}
  {F=30\N\,,\qquad \beta=60^\circ}
  {គណនាកម្លាំងបំបែកផ្ដេក $F_x$}
  {$26\N$}{$15\N$}{$30\N$}{$52\N$}{1}

% ================= 15 =================
\sol{កម្លាំងទាញទៅមុខ គឺកម្លាំងបំបែកផ្ដេក ៖\par
     $F_x=F\cos60^\circ=40\times0.5=20\N$}
\problem{មធ្យម}
  {សិស្សម្នាក់ទាញរទេះដោយខ្សែ ដែលធ្វើមុំ $\alpha$ ជាមួយផ្ទៃដីផ្ដេក ។ គេឲ្យ}
  {F=40\N\,,\qquad \alpha=60^\circ}
  {រកកម្លាំងដែលទាញរទេះទៅមុខ (កម្លាំងបំបែកផ្ដេក)}
  {$34.6\N$}{$40\N$}{$80\N$}{$20\N$}{4}

% ================= 16 =================
\sol{$P=mg=5\times10=50\N$\par
     $P_{\parallel}=P\sin30^\circ=50\times0.5=25\N$}
\problem{មធ្យម}
  {ប្រអប់ម៉ាស $m$ ស្ថិតលើប្លង់ទេរ ដែលមានមុំទេរ $\alpha$ ។ ($g=10\,\mathrm{N/kg}$) ។ គេឲ្យ}
  {m=5\kg\,,\qquad \alpha=30^\circ}
  {គណនាកម្លាំងបំបែកនៃទម្ងន់ ស្របនឹងប្លង់ទេរ $P_{\parallel}$}
  {$43.3\N$}{$25\N$}{$50\N$}{$5\N$}{2}

% ================= 17 =================
\sol{$P=mg=5\times10=50\N$\par
     $P_{\perp}=P\cos30^\circ=50\times0.866\approx43.3\N$}
\problem{មធ្យម}
  {ប្រអប់ម៉ាស $m$ ស្ថិតលើប្លង់ទេរ ដែលមានមុំទេរ $\alpha$ ។ ($g=10\,\mathrm{N/kg}$) ។ គេឲ្យ}
  {m=5\kg\,,\qquad \alpha=30^\circ}
  {គណនាកម្លាំងបំបែកនៃទម្ងន់ កែងនឹងប្លង់ទេរ $P_{\perp}$}
  {$25\N$}{$50\N$}{$43.3\N$}{$86.6\N$}{3}

% ================= 18 =================
\sol{$P=mg=2\times10=20\N$\par
     $P_{\parallel}=P\sin37^\circ=20\times0.6=12\N$}
\problem{មធ្យម}
  {ប្រអប់ម៉ាស $m$ ស្ថិតលើប្លង់ទេរ ដែលមានមុំទេរ $\alpha$ ។ ($g=10\,\mathrm{N/kg}$ ,
   $\sin37^\circ=0.6$ , $\cos37^\circ=0.8$) ។ គេឲ្យ}
  {m=2\kg\,,\qquad \alpha=37^\circ}
  {គណនាកម្លាំងបំបែកនៃទម្ងន់ ស្របនឹងប្លង់ទេរ $P_{\parallel}$}
  {$12\N$}{$16\N$}{$20\N$}{$6\N$}{1}

% ================= 19 =================
\sol{$P_{\parallel}$ និង $P_{\perp}$ កែងគ្នា ដូច្នេះ\par
     $P=\sqrt{P_{\parallel}^2+P_{\perp}^2}=\sqrt{30^2+40^2}=\sqrt{2500}=50\N$}
\problem{ពិបាក}
  {ទម្ងន់ $\vec P$ របស់វត្ថុមួយលើប្លង់ទេរ ត្រូវបានបំបែកជាកម្លាំងពីរ ។ គេឲ្យ}
  {P_{\parallel}=30\N\,,\qquad P_{\perp}=40\N}
  {រកទម្ងន់ $P$ របស់វត្ថុ}
  {$70\N$}{$10\N$}{$35\N$}{$50\N$}{4}

% ================= 20 =================
\sol{ដើម្បីទប់ប្រអប់ $F$ ត្រូវស្មើនឹង $P_{\parallel}$ ៖\par
     $F=P_{\parallel}=mg\sin30^\circ=4\times10\times0.5=20\N$}
\problem{ពិបាក}
  {ប្រអប់ម៉ាស $m$ ស្ថិតលើប្លង់ទេររលោង (គ្មានកកិត) ដែលមានមុំទេរ $\alpha$ ។
   ដើម្បីទប់ប្រអប់កុំឲ្យរអិល គេអនុវត្តកម្លាំង $\vec F$ ស្របនឹងប្លង់ ។
   ($g=10\,\mathrm{N/kg}$) ។ គេឲ្យ}
  {m=4\kg\,,\qquad \alpha=30^\circ}
  {រកអាំងតង់ស៊ីតេនៃកម្លាំង $F$}
  {$40\N$}{$20\N$}{$34.6\N$}{$10\N$}{2}

\end{document}
